% CAPITOLO 8 — VIDEO E LINK
\chapter{Video e Link}
\label{cap:links}

Il progetto è disponibile pubblicamente e la qualità è monitorata con
analisi statica. I seguenti riferimenti completano la documentazione.

\section{Repository GitHub}
Il codice sorgente completo (applicazione JavaFX/CLI, test JUnit 5, script
di build e documentazione) è disponibile sul repository pubblico:
\begin{center}
\url{https://github.com/Damiano-o/ciboamico.git}
\end{center}

\section{SonarCloud}
Il progetto è analizzato con \textbf{SonarCloud} (organizzazione
\texttt{ciboamico-ispw}); il pannello pubblico riporta il \emph{Quality
Gate}, le metriche di duplicazione, i code smells e la copertura importata da
JaCoCo. Il link di riferimento sarà:
\begin{center}
\url{https://sonarcloud.io/organizations/ciboamico-ispw}
\end{center}
La configurazione degli \emph{excludes} dei componenti grafici è dichiarata
nel \texttt{pom.xml} (sezione \texttt{sonar}) per non falsare le metriche di
copertura.

\section{Video dimostrativo}
\begin{quote}
\emph{Nota di consegna:} il video dimostrativo (\texttt{video-ubergata.mp4}) è
disponibile nel repository o su piattaforma di condivisione; il link sarà
inserito a cura dello studente prima della consegna.
\end{quote}